% --- Archon physics formalization source begin ---
% archon:physics
% archon:covers PhyXMiniProblems/problem_phyx_mini_0189.lean
% archon:source-report reports/phyx_mini/problem_phyx_mini_0189.source.json

\chapter{Physics problem phyx\_mini\_0189}
\label{ch:PhyXMiniProblems_problem_phyx_mini_0189}

\paragraph{Problem source.}
\#\# Physical scenario

A sound wave that enters the human ear sets the eardrum into oscillation, which in turn causes oscillation of the ossicles, a chain of three tiny bones in the middle ear (figure). The ossicles transmit this oscillation to the fluid (mostly water) in the inner ear; there the fluid motion disturbs hair cells that send nerve impulses to the brain with information about the sound. The area of the moving part of the eardrum is about \textbackslash{}( 43 \textbackslash{}, \textbackslash{}mathrm\{mm\textasciicircum{}2\} \textbackslash{}), and that of the stapes (the smallest of the ossicles) where it connects to the inner ear is about \textbackslash{}( 3.2 \textbackslash{}, \textbackslash{}mathrm\{mm\textasciicircum{}2\} \textbackslash{}). For the speed of sound is \textbackslash{}( 344 \textbackslash{}, \textbackslash{}mathrm\{m/s\} \textbackslash{}) and the bulk modulus is \textbackslash{}( 1.42 \textbackslash{}times 10\textasciicircum{}5 \textbackslash{}, \textbackslash{}mathrm\{Pa\} \textbackslash{}),the speed of sound in the inner ear is \textbackslash{}( 1500 \textbackslash{}, \textbackslash{}mathrm\{m/s\} \textbackslash{}).

\#\# Question

Determine the displacement amplitude of the wave in the fluid of the inner ear, in which the speed of sound is \textbackslash{}( 1500 \textbackslash{}, \textbackslash{}mathrm\{m/s\} \textbackslash{}).

\#\# Answer choices

- A: A: \$4.4 \textbackslash{}times 10\textasciicircum{}\{-11\}\textbackslash{}

- B: B: \$4.4 \textbackslash{}times 10\textasciicircum{}\{-10\}\textbackslash{}

- C: C: \$4.4 \textbackslash{}times 10\textasciicircum{}\{-12\}\textbackslash{}

- D: D: \$4.0 \textbackslash{}times 10\textasciicircum{}\{-11\}\textbackslash{}

\#\# Figure caption (auxiliary; use the image as primary evidence)

The image is a labeled diagram of the human ear, showing its internal structure.

1. **Outer Ear**:
   - **Auditory Canal**: This is seen leading to the eardrum.
   - **Eardrum**: Positioned at the end of the auditory canal, separating the outer ear from the middle ear.

2. **Middle Ear**:
   - **Ossicles (middle ear bones)**: Three small bones labeled:
     - **Malleus (anvil)**
     - **Incus (hammer)**
     - **Stapes (stirrup)**

3. **Inner Ear**:
   - **Cochlea of inner ear**: Shown as a spiral-shaped structure, which is part of the hearing process.

The diagram provides a visual representation of the ear's anatomy, with lines pointing from the labels to their corresponding structures.

\#\# Dataset metadata

Sound; Physical Model Grounding Reasoning, Multi-Formula Reasoning

\paragraph{Recorded answer/context.}
A: A: \$4.4 \textbackslash{}times 10\textasciicircum{}\{-11\}\textbackslash{}

\paragraph{Figure/image path.}
/root/proposal\_for\_physic/science-mango/phyx\_mini\_run/phyx\_data/test\_image/189.png

\paragraph{Formalization target.}
create a compiling Lean file with sorry bodies at `PhyXMiniProblems/problem\_phyx\_mini\_0189.lean`. The Lean declarations must preserve the physical quantities, dimensions or dimensional roles, figure labels, governing-law hypotheses, and final relation expressed by this problem.
Use Mathlib/Physlib names found through LeanExplore where available. If a physics API is missing, introduce faithful local abstractions rather than scalar placeholder aliases.

\begin{theorem}[Physics formalization target]
\label{thm:physics:phyx_mini_0189:target}
\lean{PhyXMiniProblems.ProblemPhyXMini0189.problem_phyx_mini_0189}
\uses{def:physics:phyx-mini-0189:phyxminiproblems-problemphyxmini0189-lengthinmeters, def:physics:phyx-mini-0189:phyxminiproblems-problemphyxmini0189-pressureinpascals, def:physics:phyx-mini-0189:phyxminiproblems-problemphyxmini0189-eartransmissionsetup, def:physics:phyx-mini-0189:phyxminiproblems-problemphyxmini0189-matchesproblemreadouts, def:physics:phyx-mini-0189:phyxminiproblems-problemphyxmini0189-hasphysicalacousticdata, def:physics:phyx-mini-0189:phyxminiproblems-problemphyxmini0189-obeysidealeartransmissionlaws}
The assigned autoformalize agent should translate this physics problem into Lean declarations in the covered file, with theorem and lemma proof bodies written as `by sorry`.
\end{theorem}
\begin{proof}
This is an autoformalization task, not a proof task. Produce statements that can later be proved by the physics prover without weakening the physical model.
\end{proof}

% --- Archon declaration topology begin ---
\section{Declaration topology}
\begin{definition}[Length Quantity]
  \label{def:physics:phyx-mini-0189:phyxminiproblems-problemphyxmini0189-lengthquantity}
  \lean{PhyXMiniProblems.ProblemPhyXMini0189.LengthQuantity}
  A nonnegative physical length, used here for a displacement amplitude
\end{definition}
\begin{proof}
  This is the declared model definition of Length Quantity.
\end{proof}
\begin{definition}[Frequency Quantity]
  \label{def:physics:phyx-mini-0189:phyxminiproblems-problemphyxmini0189-frequencyquantity}
  \lean{PhyXMiniProblems.ProblemPhyXMini0189.FrequencyQuantity}
  A nonnegative physical frequency, with dimension inverse time
\end{definition}
\begin{proof}
  This is the declared model definition of Frequency Quantity.
\end{proof}
\begin{definition}[Mass Density Quantity]
  \label{def:physics:phyx-mini-0189:phyxminiproblems-problemphyxmini0189-massdensityquantity}
  \lean{PhyXMiniProblems.ProblemPhyXMini0189.MassDensityQuantity}
  A nonnegative mass density, with dimension M L⁻³
\end{definition}
\begin{proof}
  This is the declared model definition of Mass Density Quantity.
\end{proof}
\begin{definition}[Acoustic Intensity Quantity]
  \label{def:physics:phyx-mini-0189:phyxminiproblems-problemphyxmini0189-acousticintensityquantity}
  \lean{PhyXMiniProblems.ProblemPhyXMini0189.AcousticIntensityQuantity}
  A nonnegative acoustic intensity, with SI unit W / m² and dimension M T⁻³
\end{definition}
\begin{proof}
  This is the declared model definition of Acoustic Intensity Quantity.
\end{proof}
\begin{definition}[square Millimeter Unit Choices]
  \label{def:physics:phyx-mini-0189:phyxminiproblems-problemphyxmini0189-squaremillimeterunitchoices}
  \lean{PhyXMiniProblems.ProblemPhyXMini0189.squareMillimeterUnitChoices}
  Unit choices whose length unit is the millimeter. Areas evaluated in these choices therefore have square-millimeter readouts
\end{definition}
\begin{proof}
  This is the declared model definition of square Millimeter Unit Choices.
\end{proof}
\begin{definition}[length In Meters]
  \label{def:physics:phyx-mini-0189:phyxminiproblems-problemphyxmini0189-lengthinmeters}
  \lean{PhyXMiniProblems.ProblemPhyXMini0189.lengthInMeters}
  \uses{def:physics:phyx-mini-0189:phyxminiproblems-problemphyxmini0189-lengthquantity}
  Meter readout of a physical displacement amplitude
\end{definition}
\begin{proof}
  This is the declared model definition of length In Meters.
\end{proof}
\begin{definition}[frequency In Hertz]
  \label{def:physics:phyx-mini-0189:phyxminiproblems-problemphyxmini0189-frequencyinhertz}
  \lean{PhyXMiniProblems.ProblemPhyXMini0189.frequencyInHertz}
  \uses{def:physics:phyx-mini-0189:phyxminiproblems-problemphyxmini0189-frequencyquantity}
  Hertz readout of a physical frequency
\end{definition}
\begin{proof}
  This is the declared model definition of frequency In Hertz.
\end{proof}
\begin{definition}[area In Square Meters]
  \label{def:physics:phyx-mini-0189:phyxminiproblems-problemphyxmini0189-areainsquaremeters}
  \lean{PhyXMiniProblems.ProblemPhyXMini0189.areaInSquareMeters}
  Square-meter readout of a physical area
\end{definition}
\begin{proof}
  This is the declared model definition of area In Square Meters.
\end{proof}
\begin{definition}[area In Square Millimeters]
  \label{def:physics:phyx-mini-0189:phyxminiproblems-problemphyxmini0189-areainsquaremillimeters}
  \lean{PhyXMiniProblems.ProblemPhyXMini0189.areaInSquareMillimeters}
  \uses{def:physics:phyx-mini-0189:phyxminiproblems-problemphyxmini0189-squaremillimeterunitchoices}
  Square-millimeter readout of a physical area
\end{definition}
\begin{proof}
  This is the declared model definition of area In Square Millimeters.
\end{proof}
\begin{definition}[speed In Meters Per Second]
  \label{def:physics:phyx-mini-0189:phyxminiproblems-problemphyxmini0189-speedinmeterspersecond}
  \lean{PhyXMiniProblems.ProblemPhyXMini0189.speedInMetersPerSecond}
  Meter-per-second readout of a physical propagation speed
\end{definition}
\begin{proof}
  This is the declared model definition of speed In Meters Per Second.
\end{proof}
\begin{definition}[pressure In Pascals]
  \label{def:physics:phyx-mini-0189:phyxminiproblems-problemphyxmini0189-pressureinpascals}
  \lean{PhyXMiniProblems.ProblemPhyXMini0189.pressureInPascals}
  Pascal readout of a bulk modulus. Bulk modulus has the physical dimension of pressure
\end{definition}
\begin{proof}
  This is the declared model definition of pressure In Pascals.
\end{proof}
\begin{definition}[density In Kilograms Per Cubic Meter]
  \label{def:physics:phyx-mini-0189:phyxminiproblems-problemphyxmini0189-densityinkilogramspercubicmeter}
  \lean{PhyXMiniProblems.ProblemPhyXMini0189.densityInKilogramsPerCubicMeter}
  \uses{def:physics:phyx-mini-0189:phyxminiproblems-problemphyxmini0189-massdensityquantity}
  Kilogram-per-cubic-meter readout of a mass density
\end{definition}
\begin{proof}
  This is the declared model definition of density In Kilograms Per Cubic Meter.
\end{proof}
\begin{definition}[intensity In Watts Per Square Meter]
  \label{def:physics:phyx-mini-0189:phyxminiproblems-problemphyxmini0189-intensityinwattspersquaremeter}
  \lean{PhyXMiniProblems.ProblemPhyXMini0189.intensityInWattsPerSquareMeter}
  \uses{def:physics:phyx-mini-0189:phyxminiproblems-problemphyxmini0189-acousticintensityquantity}
  Watt-per-square-meter readout of an acoustic intensity
\end{definition}
\begin{proof}
  This is the declared model definition of intensity In Watts Per Square Meter.
\end{proof}
\begin{definition}[Ear Part]
  \label{def:physics:phyx-mini-0189:phyxminiproblems-problemphyxmini0189-earpart}
  \lean{PhyXMiniProblems.ProblemPhyXMini0189.EarPart}
  Anatomical labels visible in the supplied ear diagram. The parenthetical names deliberately reproduce the labels in that figure
\end{definition}
\begin{proof}
  This is the declared model definition of Ear Part.
\end{proof}
\begin{definition}[Ear Region]
  \label{def:physics:phyx-mini-0189:phyxminiproblems-problemphyxmini0189-earregion}
  \lean{PhyXMiniProblems.ProblemPhyXMini0189.EarRegion}
  The three anatomical regions distinguished by the diagram and prose
\end{definition}
\begin{proof}
  This is the declared model definition of Ear Region.
\end{proof}
\begin{definition}[Ear Diagram Location]
  \label{def:physics:phyx-mini-0189:phyxminiproblems-problemphyxmini0189-eardiagramlocation}
  \lean{PhyXMiniProblems.ProblemPhyXMini0189.EarDiagramLocation}
  \uses{def:physics:phyx-mini-0189:phyxminiproblems-problemphyxmini0189-earregion}
  A labeled part is either within one anatomical region or forms the boundary between two regions
\end{definition}
\begin{proof}
  This is the declared model definition of Ear Diagram Location.
\end{proof}
\begin{definition}[diagram Location]
  \label{def:physics:phyx-mini-0189:phyxminiproblems-problemphyxmini0189-earpart-diagramlocation}
  \lean{PhyXMiniProblems.ProblemPhyXMini0189.EarPart.diagramLocation}
  \uses{def:physics:phyx-mini-0189:phyxminiproblems-problemphyxmini0189-earpart, def:physics:phyx-mini-0189:phyxminiproblems-problemphyxmini0189-eardiagramlocation}
  Location of each labeled part in the problem's simplified figure. The eardrum is recorded as separating the outer and middle ear, as stated in the caption, rather than being assigned arbitrarily to one side
\end{definition}
\begin{proof}
  This is the declared model definition of diagram Location.
\end{proof}
\begin{definition}[Directly Transmits Oscillation]
  \label{def:physics:phyx-mini-0189:phyxminiproblems-problemphyxmini0189-directlytransmitsoscillation}
  \lean{PhyXMiniProblems.ProblemPhyXMini0189.DirectlyTransmitsOscillation}
  \uses{def:physics:phyx-mini-0189:phyxminiproblems-problemphyxmini0189-earpart}
  Successive oscillation-transmission links described by the figure and scenario. This records anatomy rather than a quantitative transfer law
\end{definition}
\begin{proof}
  This is the declared model definition of Directly Transmits Oscillation.
\end{proof}
\begin{definition}[Ear Medium Label]
  \label{def:physics:phyx-mini-0189:phyxminiproblems-problemphyxmini0189-earmediumlabel}
  \lean{PhyXMiniProblems.ProblemPhyXMini0189.EarMediumLabel}
  The two propagation media relevant to the numerical acoustics problem
\end{definition}
\begin{proof}
  This is the declared model definition of Ear Medium Label.
\end{proof}
\begin{definition}[Acoustic Medium]
  \label{def:physics:phyx-mini-0189:phyxminiproblems-problemphyxmini0189-acousticmedium}
  \lean{PhyXMiniProblems.ProblemPhyXMini0189.AcousticMedium}
  \uses{def:physics:phyx-mini-0189:phyxminiproblems-problemphyxmini0189-massdensityquantity, def:physics:phyx-mini-0189:phyxminiproblems-problemphyxmini0189-earmediumlabel}
  Physical material data for a homogeneous acoustic medium
\end{definition}
\begin{proof}
  This is the declared model definition of Acoustic Medium.
\end{proof}
\begin{definition}[Acoustic Wave]
  \label{def:physics:phyx-mini-0189:phyxminiproblems-problemphyxmini0189-acousticwave}
  \lean{PhyXMiniProblems.ProblemPhyXMini0189.AcousticWave}
  \uses{def:physics:phyx-mini-0189:phyxminiproblems-problemphyxmini0189-lengthquantity, def:physics:phyx-mini-0189:phyxminiproblems-problemphyxmini0189-frequencyquantity, def:physics:phyx-mini-0189:phyxminiproblems-problemphyxmini0189-acousticintensityquantity}
  Scalar observables of a monochromatic acoustic wave, retained as dimensionful quantities independent of a choice of units
\end{definition}
\begin{proof}
  This is the declared model definition of Acoustic Wave.
\end{proof}
\begin{definition}[Ear Transmission Setup]
  \label{def:physics:phyx-mini-0189:phyxminiproblems-problemphyxmini0189-eartransmissionsetup}
  \lean{PhyXMiniProblems.ProblemPhyXMini0189.EarTransmissionSetup}
  \uses{def:physics:phyx-mini-0189:phyxminiproblems-problemphyxmini0189-acousticmedium, def:physics:phyx-mini-0189:phyxminiproblems-problemphyxmini0189-acousticwave}
  All physical quantities in the idealized ear-transmission experiment
\end{definition}
\begin{proof}
  This is the declared model definition of Ear Transmission Setup.
\end{proof}
\begin{definition}[Matches Problem Readouts]
  \label{def:physics:phyx-mini-0189:phyxminiproblems-problemphyxmini0189-matchesproblemreadouts}
  \lean{PhyXMiniProblems.ProblemPhyXMini0189.MatchesProblemReadouts}
  \uses{def:physics:phyx-mini-0189:phyxminiproblems-problemphyxmini0189-areainsquaremillimeters, def:physics:phyx-mini-0189:phyxminiproblems-problemphyxmini0189-speedinmeterspersecond, def:physics:phyx-mini-0189:phyxminiproblems-problemphyxmini0189-pressureinpascals, def:physics:phyx-mini-0189:phyxminiproblems-problemphyxmini0189-eartransmissionsetup}
  Numerical data explicitly printed in the problem
\end{definition}
\begin{proof}
  This is the declared model definition of Matches Problem Readouts.
\end{proof}
\begin{definition}[Has Physical Acoustic Data]
  \label{def:physics:phyx-mini-0189:phyxminiproblems-problemphyxmini0189-hasphysicalacousticdata}
  \lean{PhyXMiniProblems.ProblemPhyXMini0189.HasPhysicalAcousticData}
  \uses{def:physics:phyx-mini-0189:phyxminiproblems-problemphyxmini0189-frequencyinhertz, def:physics:phyx-mini-0189:phyxminiproblems-problemphyxmini0189-areainsquaremeters, def:physics:phyx-mini-0189:phyxminiproblems-problemphyxmini0189-speedinmeterspersecond, def:physics:phyx-mini-0189:phyxminiproblems-problemphyxmini0189-pressureinpascals, def:physics:phyx-mini-0189:phyxminiproblems-problemphyxmini0189-densityinkilogramspercubicmeter, def:physics:phyx-mini-0189:phyxminiproblems-problemphyxmini0189-eartransmissionsetup}
  Positivity and nondegeneracy conditions for the acoustic model
\end{definition}
\begin{proof}
  This is the declared model definition of Has Physical Acoustic Data.
\end{proof}
\begin{definition}[angular Frequency In Radians Per Second]
  \label{def:physics:phyx-mini-0189:phyxminiproblems-problemphyxmini0189-angularfrequencyinradianspersecond}
  \lean{PhyXMiniProblems.ProblemPhyXMini0189.angularFrequencyInRadiansPerSecond}
  \uses{def:physics:phyx-mini-0189:phyxminiproblems-problemphyxmini0189-frequencyinhertz, def:physics:phyx-mini-0189:phyxminiproblems-problemphyxmini0189-acousticwave}
  Angular-frequency readout corresponding to a wave's ordinary frequency
\end{definition}
\begin{proof}
  This is the declared model definition of angular Frequency In Radians Per Second.
\end{proof}
\begin{definition}[Obeys Ideal Ear Transmission Laws]
  \label{def:physics:phyx-mini-0189:phyxminiproblems-problemphyxmini0189-obeysidealeartransmissionlaws}
  \lean{PhyXMiniProblems.ProblemPhyXMini0189.ObeysIdealEarTransmissionLaws}
  \uses{def:physics:phyx-mini-0189:phyxminiproblems-problemphyxmini0189-lengthinmeters, def:physics:phyx-mini-0189:phyxminiproblems-problemphyxmini0189-frequencyinhertz, def:physics:phyx-mini-0189:phyxminiproblems-problemphyxmini0189-areainsquaremeters, def:physics:phyx-mini-0189:phyxminiproblems-problemphyxmini0189-speedinmeterspersecond, def:physics:phyx-mini-0189:phyxminiproblems-problemphyxmini0189-pressureinpascals, def:physics:phyx-mini-0189:phyxminiproblems-problemphyxmini0189-densityinkilogramspercubicmeter, def:physics:phyx-mini-0189:phyxminiproblems-problemphyxmini0189-intensityinwattspersquaremeter, def:physics:phyx-mini-0189:phyxminiproblems-problemphyxmini0189-eartransmissionsetup, def:physics:phyx-mini-0189:phyxminiproblems-problemphyxmini0189-angularfrequencyinradianspersecond}
  Governing linear-acoustics and ideal lossless-transmission laws. The first two fields are B = ρ c². The intensity fields state I = B ω² s² / (2c) for a plane sound wave. The final field is conservation of acoustic power through the eardrum and stapes areas. None of these fields fixes the requested inner-ear displacement to a numeric answer
\end{definition}
\begin{proof}
  This is the declared model definition of Obeys Ideal Ear Transmission Laws.
\end{proof}
\begin{definition}[Answer Choice]
  \label{def:physics:phyx-mini-0189:phyxminiproblems-problemphyxmini0189-answerchoice}
  \lean{PhyXMiniProblems.ProblemPhyXMini0189.AnswerChoice}
  Labels of the four multiple-choice answers
\end{definition}
\begin{proof}
  This is the declared model definition of Answer Choice.
\end{proof}
\begin{definition}[answer Displacement In Meters]
  \label{def:physics:phyx-mini-0189:phyxminiproblems-problemphyxmini0189-answerdisplacementinmeters}
  \lean{PhyXMiniProblems.ProblemPhyXMini0189.answerDisplacementInMeters}
  \uses{def:physics:phyx-mini-0189:phyxminiproblems-problemphyxmini0189-answerchoice}
  Displacement-amplitude readouts in meters printed in the choices
\end{definition}
\begin{proof}
  This is the declared model definition of answer Displacement In Meters.
\end{proof}
\begin{definition}[recorded Answer Choice]
  \label{def:physics:phyx-mini-0189:phyxminiproblems-problemphyxmini0189-recordedanswerchoice}
  \lean{PhyXMiniProblems.ProblemPhyXMini0189.recordedAnswerChoice}
  \uses{def:physics:phyx-mini-0189:phyxminiproblems-problemphyxmini0189-answerchoice}
  Answer A is the answer recorded by the dataset. It is retained as source metadata and is not asserted as a consequence of the printed physical data
\end{definition}
\begin{proof}
  This is the declared model definition of recorded Answer Choice.
\end{proof}
% --- Archon declaration topology end ---
% --- Archon physics formalization source end ---
